\chapter{Noticias relacionadas}

A continuación se incluyen tres fichas, una por CVE, recogiendo una noticia o publicación técnica relacionada con cada una de las vulnerabilidades analizadas.

\begin{tcolorbox}[breakable, colback=lightblue, colframe=darkblue, title=\textbf{Ficha de noticia 1}]
\textbf{CVE a la que se refiere:} CVE-2018-6350 (WhatsApp).\\
\textbf{Fecha de la noticia:} 14 de mayo de 2019.\\
\textbf{Título de la noticia:} \textit{``A WhatsApp vulnerability allowed spyware to be installed via voice calls''}.\\
\textbf{Enlace a la noticia:} \url{https://www.theguardian.com/technology/2019/may/14/whatsapp-urges-users-to-upgrade-after-discovering-spyware-vulnerability}\\
\textbf{Breve resumen:} La noticia recoge el descubrimiento, por parte de \textit{Facebook} (entonces dueña de WhatsApp) y del \textit{Citizen Lab} de la Universidad de Toronto, de una vulnerabilidad crítica en el subsistema de llamadas VoIP de WhatsApp ---catalogada con posterioridad como CVE-2019-3568--- que permitía instalar el \textit{spyware} Pegasus de NSO Group enviando una simple llamada al número de la víctima, sin necesidad de que ésta descolgara. Aunque el identificador concreto de esa campaña no es el mismo que el de CVE-2018-6350, ambas vulnerabilidades pertenecen a la misma familia técnica: errores de corrupción de memoria en el procesamiento de tramas RTP entrantes durante la negociación de la llamada, sin interacción del usuario. La noticia es relevante porque ilustra el impacto operativo real que pueden tener fallos como CVE-2018-6350 cuando son ``hermanos mayores'' o predecesores conceptuales de una campaña activa: el mismo tipo de vector se utilizó contra periodistas, abogados de derechos humanos y disidentes políticos, y obligó a WhatsApp a forzar una actualización global en cuestión de horas. Sirve, por tanto, como recordatorio de que las vulnerabilidades del parser de protocolos de mensajería instantánea no son una curiosidad académica, sino vector real de espionaje dirigido.
\end{tcolorbox}

\begin{tcolorbox}[breakable, colback=lightblue, colframe=darkblue, title=\textbf{Ficha de noticia 2}]
\textbf{CVE a la que se refiere:} CVE-2025-53766 (Microsoft GDI+).\\
\textbf{Fecha de la noticia:} 13 de agosto de 2025.\\
\textbf{Título de la noticia:} \textit{``Microsoft August 2025 Patch Tuesday fixes 107 flaws, including a critical wormable GDI+ RCE''}.\\
\textbf{Enlace a la noticia:} \url{https://www.bleepingcomputer.com/news/microsoft/microsoft-august-2025-patch-tuesday/}\\
\textbf{Breve resumen:} El boletín mensual de Microsoft de agosto de 2025 fue uno de los más voluminosos del año, con más de cien vulnerabilidades corregidas en un solo ciclo. Entre ellas, \textit{BleepingComputer} destaca específicamente CVE-2025-53766 por dos motivos: su puntuación de 9.8 CRITICAL y, sobre todo, el hecho de que el componente afectado (GDI+) es invocado de manera transparente por una larguísima lista de procesos del sistema, incluidos servicios que decodifican imágenes en cadenas de procesamiento automáticas (motores de \textit{indexing}, vistas previas, miniaturas, pasarelas de correo). El artículo subraya que esta naturaleza ``en cadena'' eleva el riesgo de transformar la vulnerabilidad en un vector \textit{wormable} si se descubriera un servicio sistémicamente accesible que la utilice, aunque a fecha de publicación de la noticia Microsoft no había reportado evidencia de explotación activa. La pieza también detalla las recomendaciones operativas habituales: aplicar el parche acumulativo, revisar las políticas de procesamiento automatizado de imágenes en servidores y considerar reglas de EDR para detectar comportamientos anómalos en procesos que carguen \texttt{gdiplus.dll}. Sirve como ejemplo claro del modelo de divulgación coordinada de Microsoft (\textit{Patch Tuesday}) y de la importancia de monitorizar los boletines mensuales.
\end{tcolorbox}

\begin{tcolorbox}[breakable, colback=lightblue, colframe=darkblue, title=\textbf{Ficha de noticia 3}]
\textbf{CVE a la que se refiere:} CVE-2023-51582 (Voltronic Power ViewPower Linux).\\
\textbf{Fecha de la noticia:} 22 de diciembre de 2023.\\
\textbf{Título de la noticia:} \textit{``ZDI-23-1887: Voltronic Power ViewPower Pro LinuxMonitorConsole Exposed Dangerous Method Remote Code Execution Vulnerability''}.\\
\textbf{Enlace a la noticia:} \url{https://www.zerodayinitiative.com/advisories/ZDI-23-1887/}\\
\textbf{Breve resumen:} El aviso oficial de \textit{Zero Day Initiative} que dio lugar a CVE-2023-51582 detalla, sin entrar en código de explotación funcional, el fallo descubierto en el componente \texttt{LinuxMonitorConsole} del agente Voltronic Power ViewPower (también comercializado bajo otras marcas blancas en el sector de SAI). El aviso describe cómo el servicio expone un método remoto que debería estar reservado a operaciones administrativas internas pero, por una configuración insuficiente del control de acceso, queda alcanzable por cualquier cliente capaz de conectarse al puerto de gestión, lo que se traduce en una ejecución remota de código en el contexto del usuario bajo el que corre el agente ---típicamente con privilegios elevados sobre la máquina anfitriona porque controla el SAI físico. Además del detalle técnico, el aviso es relevante porque pone el foco sobre una categoría de \textit{software de gestión de infraestructura} que recibe mucha menos atención mediática que las grandes plataformas (Windows, macOS, Android), pero que está omnipresente en CPDs e instalaciones industriales. La pieza concluye con la recomendación de actualizar a la última versión del agente y, mientras tanto, restringir mediante \textit{firewall} el acceso al servicio.
\end{tcolorbox}
